% Figura 9.1: Modigliani-Miller Proposición II sin impuestos frente a la Tesis Tradicional (Tema 7 UGR)
\begin{tikzpicture}
\begin{groupplot}[
    group style={group size=2 by 1, horizontal sep=1.55cm},
    width=6.35cm, height=6.0cm,
    xlabel={Ratio de endeudamiento $D/E$ (valor de mercado)},
    ylabel={Rentabilidad exigida / Coste (\%)},
    xmin=0, xmax=2.2,
    ymin=0, ymax=24,
    xtick={0, 0.5, 1.0, 1.5, 2.0},
    ytick={0, 4, 8, 12, 16, 20},
    yticklabels={{0\%}, {4\%}, {8\%}, {12\%}, {16\%}, {20\%}},
    grid=major, grid style={dotted, color=dfinmuted!45},
    axis line style={color=dfinslate, line width=0.9pt},
    tick label style={font=\scriptsize\sffamily, color=dfinslate},
    label style={font=\footnotesize\sffamily\bfseries, color=dfinnavy},
    title style={font=\footnotesize\sffamily\bfseries, color=dfinnavy},
    legend style={at={(0.03,0.97)}, anchor=north west, font=\tiny\sffamily,
                  fill=white, draw=dfinslate!35, rounded corners=2pt}
]

% Panel izquierdo: Modigliani-Miller sin impuestos (rA = 10% constante, rD = 4% con ligero repunte en riesgo extremo)
\nextgroupplot[title={MM II (sin impuestos): WACC constante}]
% rA = 10% constante
\addplot[domain=0:2.1, samples=2, color=dfinnavy, line width=1.65pt] {10};
\addlegendentry{Coste activos $r_A = \text{CMPC} = 10\%$}

% rD crece suavemente al pasar D/E > 1.2 (deuda con riesgo), y rE compensa curvándose ligeramente hacia abajo
% Definimos rD(x) = 4 + 1.8*(max(0, x-1.1))^1.5 ; rE(x) = 10 + (10 - rD(x))*x
\addplot[domain=0:2.1, samples=55, color=dfincopper, line width=1.45pt]
    { 10 + (10 - (4 + (x>1.1 ? 1.8*(x-1.1)^1.5 : 0))) * x };
\addlegendentry{Coste acciones $r_E$}

\addplot[domain=0:2.1, samples=55, color=dfinemerald, line width=1.35pt, dashed]
    { 4 + (x>1.1 ? 1.8*(x-1.1)^1.5 : 0) };
\addlegendentry{Coste deuda $r_D$}

% Panel derecho: Tesis Tradicional (WACC en U con mínimo artificial)
\nextgroupplot[title={Tesis Tradicional: mínimo en $(D/E)^*$}]
% rE apenas sube al principio y luego se dispara
\addplot[domain=0:2.1, samples=50, color=dfincopper, line width=1.45pt]
    { 10 + 0.6*x + 2.4*x^2 };
\addlegendentry{$r_E$ (sube lento al principio)}

% rD sube progresivamente desde D/E = 0.8
\addplot[domain=0:2.1, samples=50, color=dfinemerald, line width=1.35pt, dashed]
    { 4 + 2.2*x^1.6 };
\addlegendentry{$r_D$}

% WACC ponderado: (1/(1+x))*rE + (x/(1+x))*rD -> dibuja forma de U clara alrededor de x=0.7
\addplot[domain=0:2.1, samples=60, color=dfinnavy, line width=1.65pt]
    { (1/(1+x))*(10 + 0.6*x + 2.4*x^2) + (x/(1+x))*(4 + 2.2*x^1.6) };
\addlegendentry{CMPC en forma de U}

% Marcar mínimo aproximado en x = 0.65 -> (1/1.65)*(10+0.39+1.01) + (0.65/1.65)*(4+1.10) = 6.91 + 2.01 = 8.92%
\addplot[only marks, mark=*, mark size=2.6pt, color=dfinburg] coordinates {(0.65, 8.92)};
\draw[dotted, color=dfinburg, line width=0.9pt] (axis cs:0.65, 0) -- (axis cs:0.65, 8.92);
\node[font=\tiny\sffamily\bfseries, color=dfinburg, anchor=north, fill=white, inner sep=1pt]
    at (axis cs:0.75, 6.2) {Óptimo $(D/E)^*$};

\end{groupplot}
\end{tikzpicture}
